\documentclass[11pt]{article}
\usepackage{preamble}
\renewcommand{\answer}[1]{}

\begin{document}

\ladderheading{AP Statistics \textperiodcentered\ Unit 3 \textperiodcentered\ Topic 3.14 \textperiodcentered\ B: 2026-12-18 \textperiodcentered\ E: 2027-02-01}{Two Simulations for Mystery Tokens}

\focusbanner{TODAY'S PROBLEM}

Suppose a company claims 10{,}000 mystery tokens are $50\%$ red, $30\%$ teal, and $20\%$ gold. An SRS of 80 has the shown counts.\par\smallskip \textbf{Leo:} ``Mix and redraw these same 80 tokens without replacement; recompute $\chi^2$ for 1,000 trials.''\par \textbf{Amara:} ``For each trial, generate 80 independent colors with probabilities 0.50, 0.30, and 0.20, then recompute $\chi^2$.'' In Amara's 1,000 trials, 186 statistics are at least the observed value.\par
\begin{center}
\textbf{Token sample}\par\smallskip
\begin{tabular}{|l|c|c|}
\hline
\textbf{Color} & \textbf{Claim} & \textbf{Observed} \\ \hline
\textbf{Red} & $0.50$ & 32 \\ \hline
\textbf{Teal} & $0.30$ & 30 \\ \hline
\textbf{Gold} & $0.20$ & 18 \\ \hline
\end{tabular}
\end{center}
\begin{firsttakebox}{FIRST TAKE --- before the video (5 min)}
Which plan can produce a different number of red tokens on the next trial? Explain in one sentence.
\workspace{0.9in}
\end{firsttakebox}

\begin{callout}{\IconBook\ EXPECTED COUNTS AND SIMULATION (keep on the page)}{calloutgreen}
For a model probability $p_i$ and sample size $n$, the expected count is $E_i=np_i$.
The chi-square statistic adds the relative squared gaps: $\chi^2=\sum (O-E)^2/E$.
To model chance variation, generate new samples of the same size using the claimed probabilities.
Compute the statistic each time. The fraction at least as large as the observed statistic estimates
how often such a mismatch occurs if the model is true. A fraction at or below a stated cutoff is
evidence against the model; a larger fraction does not prove the model.
\end{callout}

\newpage
\textbf{Finish after the video:}\par\smallskip

\dokbadge{1}~\textbf{(a)} Find the three expected color counts from the company's percentages.\par
\workspace{0.7in}

\dokbadge{2}~\textbf{(b)} The observed statistic is $\chi^2=3.35$. Check the red contribution $(32-40)^2/40$, then compute the simulation proportion $186/1{,}000$.\par
\workspace{1.4in}

\dokbadge{3}~\textbf{(c)} In 3--4 sentences, choose a simulation and justify what varies from trial to trial. Use the simulated proportion to decide whether the result is unusual at a 5 percent cutoff, and explain why this does not prove the company's claim.\par
\workspace{2.0in}

\begin{sentenceframebox}\raggedright\textbf{Frame:} I choose \blank[1.1in] because the counts \blank[1.8in] across trials. The other plan keeps \blank[1.8in] fixed.\end{sentenceframebox}

\begin{sentenceframebox}\raggedright\textbf{Frame:} Since the simulated proportion is \blank[0.8in], the result \blank[1.8in]; this does not prove \blank[1.8in].\end{sentenceframebox}

\begin{summaryexitbox}{EXIT --- turn this in}
One thing the video changed about my first take: \hrulefill\par
\workspace{0.4in}
\end{summaryexitbox}

\end{document}
